\documentclass[a4paper,12pt]{article}
\usepackage[utf8]{inputenc}
\usepackage[T1]{fontenc}
\usepackage[hungarian]{babel}
\usepackage{graphicx}
\usepackage{geometry}
\geometry{a4paper,
		     tmargin = 30mm, 
		     lmargin = 30mm,
		     rmargin = 30mm,
		     bmargin = 30mm}
\usepackage{mathtools}
\usepackage{amsmath}
\usepackage{color}
\usepackage{setspace}
\usepackage{amsmath,amssymb}
\usepackage{float}

\begin{document}

\begin{titlepage}
% Template for future notes.

	\centering
	
	\includegraphics[width=0.66\textwidth]{elte.jpg}\par\vspace{1cm}
	{\scshape\LARGE ELTE TTK \par}
	{\scshape\large Fizika szak \par}
	\vspace{1cm}
	
	{\scshape\LARGE Szakdolgozat \par}
	\vspace{1cm}
	
	{\scshape\Large Koaleszcencia modell kidolgozása transzportmodellekhez\par}
	\vspace{0.2cm}
	{\Large\itshape Olar Alex\par}
	\vspace{2cm}
	
	témavezető\par
	\vspace{0.3cm}
	{\Large Wolf György}

	\vfill

	{\large 2018 \par}
\end{titlepage}

\begin{abstract}
\par Egy nehézion ütközésben résztvevő alkotó elemek száma néhány ezerig terjed legfeljebb, így a kidolgozott néhány-test elméletek, mint a három-test problémára kidolgozott Fagyejev-egyenletek, nem alkalmazhatóak, de az alkotóelemek alacsony száma miatt még a statisztikus fizikai modellek sem használhatóak, ráadásul nem is egyensúlyi reakciókról van szó az esetek többségében.
\par A rendelkezésre álló számítási kapacitás lehetővé tette, egy-egy ilyen nemegyensúlyi reakció teljes vizsgálatát, mikroszkópikus transzport-modellek segítségével. Egy ilyen modell a Boltzmann-Uehling-Uhlenbeck elmélet (BUU), ami fázistérben leírja adott részecskék között az ütközéseket és figyelembe veszi az azok között ható kölcsönhatást, egy időfüggő, átlagtér potenciállal. Korai modellek a részecskéket szabadnak tekintették, amikor azok nem vettek részt ütközésekben.
\par Az én célom, hogy egy, a BUU-ra épülő szimulációhoz kidolgozzak egy olyan programot, ami a kölcsönható részecskéket, esetemben főként nukleonokat, klaszterezi, azaz csomósodásokat keres különböző távolság definíciók mellett (térben, impulzustérben, stb.). Ennek fontos szerepe lehet a detektor válasz meghatározásakor, 
\end{abstract}

\vfill

\tableofcontents

\newpage
%%%%%%%%%%%%%%%%%%%%%


\end{document}